% !TeX root = ../教材题库_全部条目_转写工作稿.tex
% Source: XJ-B1-C05-EX-P159-Y1414
% 第5章 / 5.1 任意角和弧度制 / 教材 p.159 / PDF 第166页
% OCR未校对：公式与文字均可能有误；未完成数学TeX转写；禁止直接发给学生
% 原题及相邻原解截图（核对用）：../题目原图/XJ-B1-C05-EX-P159-Y1414-01.webp
% 本文件只在转写完成后作为题干使用；OCR 条目尚不是可靠题干。
\begin{bankitem}{5.1 例题 例题3}{第5章 / 5.1 任意角和弧度制 / 教材 p.159 / PDF 第166页}{XJ-B1-C05-EX-P159-Y1414}{OCR未校对：公式与文字均可能有误；未完成数学TeX转写；禁止直接发给学生}
\begin{ocrdraft}
[教材p.159 / PDF第166页]
\par
0河把下列各角从度化为弧度：
\par
（1）120。；
\par
（2）25。30'．
\par
解（1)120。=120\ensuremath{\times}
\par
rad=
\par
rad；
\par
180
\par
3
\par
17丌
\par
（2）25。30'一25．5。=25．5\ensuremath{\times}
\par
rad=
\par
rad.
\par
180
\par
120
\end{ocrdraft}
\end{bankitem}
